\subsection{Statement of the flatness criterion}

Let A be a ring, \(\mathfrak{I}\) a two-sided ideal of A, M a left-module and gr(A) and gr(M) the graded ring and graded gr(A)-module associated respectively with the ring A and with the module M with the \(\mathfrak{I}\) -adic filtrations (§ 2, no. 3). We have seen (loc. cit.) that for every integer \(n \geqslant 0\) there is a surjective Z-module homomorphism

\[
\gamma_ {n} \colon (\mathfrak {I} ^ {n} / \mathfrak {I} ^ {n + 1}) \otimes_ {\mathrm{A} / \mathfrak {I}} (\mathrm{M} / \mathfrak {I} \mathrm{M}) \rightarrow \mathfrak {I} ^ {n} \mathrm{M} / \mathfrak {I} ^ {n + 1} \mathrm{M}
\]

and a graded homomorphism of degree 0 of graded gr(A)-modules

\[
\gamma_ {M} \colon \operatorname{gr} (A) \otimes_ {\mathbf {g r} _ {0} (A)} \operatorname{gr} _ {0} (M) \to \operatorname{gr} (M)
\]

whose restriction to \(\mathrm{gr}_n(\mathbf{A})\otimes_{\mathbf{gr}_0(\mathbf{A})}\mathrm{gr}_0(\mathbf{M})\) is \(\gamma_{n}\) for all \(n\) and which is therefore surjective.

THEOREM 1. Let A be a commutative ring, \(\Im\) an ideal of A and M an A-module. Consider the following properties:

\begin{enumerate}
\def\labelenumi{(\roman{enumi})}
\item
  M is a flat A-module.
\item
  \(\operatorname{Tor}_1^{\mathbf{A}}(\mathbf{N},\mathbf{M}) = 0\) for every A-module \(\mathbf{N}\) annihilated by \(\Im\) .
\item
  M/3M is a flat (A/3)-module and the canonical mapping \(\Im \otimes_{\mathbf{A}} \mathbf{M} \to \Im \mathbf{M}\) is bijective (the latter condition being equivalent to \(\operatorname{Tor}_{1}^{\mathbf{A}}(\mathbf{A}/\Im, \mathbf{M}) = 0\) by virtue of the relation \(\operatorname{Tor}_{1}^{\mathbf{A}}(\mathbf{A}, \mathbf{M}) = 0\) and the exact sequence

\[
\operatorname{Tor} _ {1} ^ {\mathbf {A}} (\mathrm{A}, \mathrm{M}) \rightarrow \operatorname{Tor} _ {1} ^ {\mathbf {A}} (\mathrm{A} / \Im , \mathrm{M}) \rightarrow \Im \otimes_ {\mathrm{A}} \mathrm{M} \rightarrow \mathrm{M}).
\]

\item
  M/3M is a flat (A/3)-module and the canonical homomorphism

\[
\gamma_ {\mathbf {M}} \colon \operatorname{gr} (\mathbf {A}) \otimes_ {\mathbf {g r} _ {0} (\mathbf {A})} \operatorname{gr} _ {0} (\mathbf {M}) \rightarrow \operatorname{gr} (\mathbf {M})
\]

is bijective (property (GR) of § 2, no. 8).

\item
  For all \(n \geqslant 1\) , \(\mathbf{M} / \Im^n\mathbf{M}\) is a flat \((\mathbf{A} / \Im^n)\) -module.
\end{enumerate}

\[
(i) \Rightarrow (i i) \Leftrightarrow (i i i) \Rightarrow (i v) \Leftrightarrow (v).
\]

If further \(\mathfrak{I}\) is nilpotent or if A is Noetherian and M is ideally Hausdorff, properties (i), (ii), (iii), (iv) and (v) are equivalent.

Remark. If A/3 is a field (as often happens in applications) the condition ``M/3M is a flat (A/3)-module'' holds automatically for every A-module M, which simplifies the statement of properties (iii) and (iv); moreover, in this case, property (v) is equivalent to saying that M/3 \(^{n}\) M is a free (A/3 \(^{n}\) )-module for every integer n ≥ 1 (Chapter II, § 3, no. 2, Corollary 2 to Proposition 5).

